\section{Visualisierungen}
\label{sec:visualisierungen}

Dieses Kapitel leitet die Darstellung vom Zielproblem her. Zunächst werden
Aufgaben ohne Festlegung auf Diagrammtypen analysiert. Daraus folgen
Anforderungen; erst anschlie{\ss}end werden die drei visuellen Kodierungen und
ihre Alternativen begründet. \Cref{fig:overview} zeigt die gemeinsame
Oberfläche in der Standardansicht.

\begin{figure}[p]
  \centering
  \includegraphics[width=\textwidth,height=0.79\textheight,keepaspectratio]{01_uebersicht.png}
  \caption{Gesamtansicht der Elm-Anwendung mit Netzwerk, Zeitreihe und
    Pixelmatrix. Alle Zahlen und SVG-Geometrien werden von Elm aus dem
    HTTP-Datensatz erzeugt.}
  \label{fig:overview}
\end{figure}

\subsection{Analyse der Anwendungsaufgaben}
\label{sec:aufgaben}

Aus der Forschungsfrage entstehen drei miteinander verbundene
Anwendungsaufgaben. Sie werden in \cref{tab:tasks} unabhängig von einer
konkreten Visualisierung formuliert.

\begin{table}[htbp]
  \centering
  \caption{Anwendungsaufgaben und dafür benötigte mentale Modelle.}
  \label{tab:tasks}
  \begin{tabularx}{\textwidth}{@{}p{0.08\textwidth}p{0.39\textwidth}X@{}}
    \toprule
    ID & Aufgabe & aufzubauendes mentales Modell \\
    \midrule
    A1 & Für eine Stunde erkennen, zwischen welchen Partnerländern und
      Deutschland physische Leistung in welcher Richtung und Grö{\ss}enordnung
      flie{\ss}t. & Deutschland als Zentrum eines gerichteten, gewichteten Netzes;
      Länderbeziehungen statt isolierter Tabellenzeilen. \\
    A2 & Zeitliche Muster und Ausnahmen in Erzeugungsmix, Preis und dem Fluss
      eines ausgewählten Länderpaars untersuchen. & Gemeinsame Zeitachse mit
      mehreren synchronen quantitativen Grö{\ss}en und klar getrennten Skalen. \\
    A3 & In vielen Länder-Stunden-Kombinationen Perioden mit ähnlicher Richtung,
      starken Beträgen oder Richtungswechseln finden und als Detailzustand
      auswählen. & Dichte zweidimensionale Übersicht mit Land und Zeit als
      orthogonalen Dimensionen. \\
    \bottomrule
  \end{tabularx}
\end{table}

A1 verlangt keine geografisch exakte Karte. Entscheidend ist die
Nachbarschaftsbeziehung zu Deutschland; Distanzen zwischen Ländern sind für den
Datensatz nicht definiert. A2 verlangt eine geordnete Zeitachse und quantitative
Ablesbarkeit. A3 verlangt Überblick vor Detail, weil \(744\times11=8184\)
physische Flusswerte nicht einzeln durchsucht werden sollen.

Übergreifend muss ein gefundener Zustand zwischen den mentalen Modellen
übersetzbar sein: Das in A3 entdeckte Pixel soll dasselbe Land im Beziehungsnetz
und dieselbe Stunde in der Zeitreihe aktivieren. Koordinierte Ansichten senken
diesen Übersetzungsaufwand. North und Shneiderman zeigen für
Overview-and-Detail-Kopplungen Leistungsgewinne gegenüber unkoordinierten
Ansichten \cite{north2000snap}.

\subsection{Anforderungen an die Visualisierungen}
\label{sec:anforderungen}

Aus den Aufgaben werden folgende Anforderungen abgeleitet, weiterhin ohne
einen bestimmten Diagrammtyp vorauszusetzen:

\begin{description}[style=nextline]
  \item[R1 -- Richtung und Betrag.] Positive und negative Flüsse müssen
    unterscheidbar sein; der Betrag soll zusätzlich quantitativ oder zumindest
    ordinal ablesbar bleiben.
  \item[R2 -- Länderidentität.] Alle elf Partner müssen gleichzeitig erkennbar
    und einzeln auswählbar sein.
  \item[R3 -- gemeinsame Zeit.] Erzeugung, Preis und Länderfluss müssen auf
    denselben Stundenschritt bezogen werden können.
  \item[R4 -- absolute und relative Werte.] Die Anwendung muss Leistung in GW,
    Preis in EUR/MWh und Erzeugungsanteile anzeigen, ohne unvereinbare Einheiten
    auf eine gemeinsame Achse zu zwingen.
  \item[R5 -- skalierbarer Überblick.] Ein voller Monat muss als Musterfläche
    betrachtbar sein; für genaues Ablesen sind kleinere Zeitfenster nötig.
  \item[R6 -- vergleichbare Kodierung.] Gleiche Farben müssen über Länder und
    Zeitfenster dieselbe Bedeutung besitzen. Neutrale Werte dürfen nicht wie
    fehlende Daten aussehen.
  \item[R7 -- verknüpfte Selektion.] Jede Auswahl muss sichtbar sein und die
    anderen Ansichten in einen konsistenten Zustand versetzen.
  \item[R8 -- semantische Transparenz.] Einheit, Vorzeichen, Quelle und
    Unsicherheiten müssen in der Oberfläche oder unmittelbar im Bericht
    erklärt werden.
\end{description}

\subsection{Präsentation der Visualisierungen}
\label{sec:praesentation}

\subsubsection{Visualisierung Eins: gerichteter radialer Netzwerkgraph}
\label{sec:vis1}

Deutschland bildet den zentralen Knoten, elf Partnerländer liegen radial um
ihn. Jede Kante ist gerichtet und gewichtet. Rot plus Pfeil nach Deutschland
steht für Import; Blau plus Pfeil vom Zentrum steht für Export. Richtung wird
damit redundant durch Geometrie und Farbe kodiert. Die Linienbreite bildet den
absoluten CBPF-Betrag ab. Ein Tooltip ergänzt den exakten physischen und den
kommerziellen Wert.

Diese Darstellung erfüllt A1 und R1--R2 direkt: Beziehungen sind wichtiger als
eine geografische Position, und alle Partner sind gleichzeitig sichtbar. Die
Längen der Kanten tragen keine Datenbedeutung. Dadurch wird vermieden, dass
eine schematische Länderposition als reale Leitungslänge interpretiert wird.
Im Feedback zum zweiten Zwischenstand wirkten Pfeilspitzen zu gro{\ss} und starke
Flüsse zu wenig differenziert. Die Marker wurden deshalb von 6 auf 4
SVG-Einheiten verkleinert; die maximal mögliche Linienbreite wurde deutlich
erhöht.

Eine Europakarte wäre eine echte Alternative: Sie unterstützt geografische
Orientierung und könnte reale Grenzlagen zeigen. Für nur elf Beziehungen zu
einem Zentrum benötigt sie jedoch mehr Fläche, und lange Landesnamen sowie
kleine Länder erschweren Auswahl und Beschriftung. Eine Tabelle wäre genauer,
verliert aber das gerichtete Beziehungsmodell. Der radiale Graph priorisiert
daher strukturellen Überblick, während Tooltips die exakte Zahl liefern.

\subsubsection{Visualisierung Zwei: gestapelte Zeitreihe mit Flussdetail}
\label{sec:vis2}

Vier gestapelte Flächen zeigen die absolute Erzeugungsleistung in GW. Die
linke Y-Achse beginnt bei null und wird aus dem sichtbaren Fenster auf einen
lesbaren Maximalwert gerundet. Ein darunterliegender, zeitlich ausgerichteter
Teilplot zeigt den CBPF-Wert des ausgewählten Partnerlands auf einer eigenen
symmetrischen GW-Skala. Eine gestrichelte Vertikale verbindet beide Teilplots.
Die Beschriftung ``+ Import nach Deutschland, -- Export aus Deutschland''
erklärt das Vorzeichen sichtbar.

Die Trennung der Skalen ist wesentlich: Erzeugungsleistungen von mehreren
Zehner-GW und bilaterale Flüsse von wenigen GW wären auf derselben Achse nicht
gleichzeitig lesbar. Die gemeinsame X-Achse erfüllt dennoch R3. Unter dem SVG
stehen für die ausgewählte Stunde Gesamtleistung, absolute Gruppenwerte,
Anteile, Preis und Länderfluss. So erfüllt die Ansicht R4, ohne den
Day-Ahead-Preis als fachlich unpassende zweite Linie auf die Leistungsachse zu
legen. Grundprinzipien zeitbezogener Datenvisualisierung sind ausführlich bei
Aigner et al. beschrieben \cite{aigner2011time}.

Als Alternative wäre ein 100-Prozent-Flächendiagramm geeignet, Anteile zu
vergleichen. Es würde aber absolute Produktionsspitzen verbergen. Mehrere
Small-Multiple-Linien wären für Einzeltechnologien präziser, erschweren jedoch
das Ablesen der Gesamtleistung und benötigen mehr vertikale Fläche. Die
gewählte Kombination zeigt den Gesamtumfang und die Zusammensetzung; die
Detailkarten kompensieren das schwierigere Ablesen mittlerer Stapelschichten.

\subsubsection{Visualisierung Drei: Pixelmatrix}
\label{sec:vis3}

Die Pixelmatrix ordnet Partnerländer zeilenweise und Stunden spaltenweise an.
Jede Zelle verwendet eine divergierende Farbskala: gesättigtes Blau bedeutet
starken Export, gesättigtes Rot starken Import, Hellgrau liegt nahe null. Eine
einzige globale Grenze wird aus allen 8184 CBPF-Werten des Monats berechnet.
Deshalb ist ein Rotton in der Frankreich-Zeile quantitativ mit demselben Ton in
der Dänemark-Zeile und in jedem Zeitfenster vergleichbar. Die Legende benennt
beide Skalenenden in GW.

Die Matrix erfüllt A3 sowie R5--R6. Pixelorientierte Verfahren nutzen die
Bildfläche effizient, weil jeder Messwert auf ein kleines farbiges Element
abgebildet wird \cite{keim2000pixel}. Im zweiten Zwischenstand trennten wei{\ss}e
Konturen die Zellen. Sie wurden entfernt, da Wei{\ss} sonst wie fehlende Daten
wirken und schmale Muster zerschneiden konnte. Der neutrale Bereich ist jetzt
bewusst hellgrau. Bei einer Auswahl markieren goldene Konturen die vollständige
Zeile und Spalte; die aktive Zelle erhält zusätzlich eine dunkle Kontur.

Elf übereinandergelegte Linien wären eine Alternative, erzeugten aber viele
Überdeckungen und machten Ländervergleiche schwierig. Small Multiples mit elf
separaten Linien böten bessere quantitative Ablesbarkeit, benötigten jedoch
deutlich mehr Fläche. Die Pixelmatrix ist als Überblick und Auswahlwerkzeug
konzipiert; exakte Werte werden nach Auswahl in Tooltip, Zeitreihe und
Detailkarten gelesen.

\subsection{Interaktion}
\label{sec:interaktion}

Die Kopplung folgt einem gemeinsamen Elm-Zustand. Ein Landklick im Netzwerk
aktualisiert die Partnerauswahl; die Matrixzeile wird hervorgehoben und der
untere Zeitreihenplot erscheint. Ein Klick auf die Zeitreihe aktualisiert die
Zeitwahl; Netzwerkwerte, vertikale Markierung und Matrixspalte wechseln
gleichzeitig. Eine Matrixzelle setzt beide Werte in einem Schritt.
Diese Operation bietet einen zusätzlichen Vorteil gegenüber der ohnehin
gemeinsamen X-Achse: Sie wählt zugleich eine von elf Länderdimensionen und eine
Stunde.

Die sichtbaren Zeitfenster 48 Stunden, sieben Tage und gesamter Monat erfüllen
R5. Vor- und Zurück-Schaltflächen verschieben kleinere Fenster ohne
Datenüberlauf. Beim Fensterwechsel bleibt die Partnerwahl erhalten; die
Zeitwahl wird auf einen gültigen Index gesetzt. ``Auswahl zurücksetzen'' löscht
den Länderfilter und setzt die Zeit an den Beginn des aktuellen Fensters.

Nicht umgesetzt wurden freies Zoomen, Mehrfachauswahl und Sortieren der
Matrixzeilen. Sie würden zusätzliche Zustände und Konflikte erzeugen, ohne für
die drei Kernaufgaben nötig zu sein. Die feste Reihenfolge erleichtert
Vergleiche zwischen Fenstern; drei wohldefinierte Zeitskalen sind für den
Monatsdatensatz verständlicher als ein unbeschrifteter Zoomfaktor.
